Predicting transcriptional responses to unseen genetic perturbations is a central task in single-cell biology, because perturbation effects depend on cellular context and spread through nonlinear regulatory interactions.
Existing predictors rely on statistical associations or unconstrained graph learning, so they rarely encode regulatory mechanisms or explain predictions in biological terms.
We propose iPerturb, a biologically informed and explainable framework that couples context-specific gene regulatory network (GRN) construction with learnable nonlinear perturbation propagation.
First, iPerturb assembles a GRN per cell line from five evidence tiers, separating a signed, directed regulatory core from undirected association evidence and deriving response-based quantities from training perturbations only.
Second, it propagates a clamped perturbation with regulator-shared Hill transfer functions, which have a provably unique steady state and run as matrix products, and integrates multi-hop propagation across tiers.
On \datanum{n_uniq_total} held-out CRISPRi targets in K562 and RPE1, iPerturb raises the centred response correlation over a linear baseline from \datanum{k562_b2_rc} to \datanum{g12_k562_prop_integrated__nnls} and from \datanum{rpe1_b2_rc} to \datanum{g12_rpe1_prop_integrated__nnls}, exceeds GEARS, CPA and biolord, and loses its gain on degree-preserving rewired networks.
The Hill operator matches exact Hill propagation to \datanum{eq_cheb16_spmm_dx} with \datanum{memratio_lo}- to \datanum{memratio_hi}-fold less memory, and its signed paths predict the direction of first-hop responses in \datanum{mech_k562_real_h1_pct}\% and \datanum{mech_rpe1_real_h1_pct}\% of responding genes.
iPerturb links accurate perturbation-response prediction to regulatory evidence that can be inspected edge by edge.
